%%═══════════════════════════════════════════════════════════════════
%% Lecture 07 — Nuclear Detectors: Resolution, Efficiency, Tracking
%% 77603 — Nuclear Physics
%% Date: 26 May 2026
%% Lecturer: Dr. Moshe Friedman
%%═══════════════════════════════════════════════════════════════════

\section{Lecture 07 — Nuclear Detectors: Resolution, Efficiency, Tracking}
\label{sec:lec07}
\date{26 May 2026}

%%───────────────────────────────────────────────────────────────────
\subsection*{Overview}
%%───────────────────────────────────────────────────────────────────
This lecture begins with a correction to the Bethe-Bloch formula from Lecture~6,
clarifies the mass-dependence of stopping power, and then develops the remaining
detector topics: the statistical origin of energy resolution, detector efficiency,
particle tracking, particle identification, time-of-flight for neutrons, and a
look at real experimental spectra.

%%───────────────────────────────────────────────────────────────────
\subsection{Correction to Lecture 6: Mass Dependence in Bethe-Bloch}
\label{subsec:lec07:bethe_correction}

\subsubsection*{Physical Motivation}
In Lecture~6 we quoted the Bethe-Bloch formula and stated that the projectile
mass does \emph{not} appear.  A question from a student prompted a more careful
discussion: where does the projectile mass enter, and why can the range of an
electron be larger than that of a proton at the same kinetic energy?

\paragraph{Why the projectile mass is absent from $-dE/dx$.}
Recall how we derived the momentum transfer to a single electron at impact
parameter $b$:
\begin{equation}
  \label{eq:lec07:delta_p}
  \Delta p_\perp = \frac{ze^2}{2\pi\epsilon_0\,b\,v}.
\end{equation}
The energy transferred to that electron is
\begin{equation}
  \label{eq:lec07:delta_E}
  \Delta E(b) = \frac{(\Delta p_\perp)^2}{2m_e} = \frac{z^2 e^4}{8\pi^2\epsilon_0^2 m_e v^2 b^2}.
\end{equation}
The \emph{electron} mass $m_e$ appears — not the projectile mass $M$.  This is
because the entire derivation was formulated in terms of \emph{velocity}: the
projectile's mass never entered the impulse integral.  The factor $1/(2m_e)$ is
the electron's kinetic energy for a given momentum kick.

\dfn[dfn:Wvalue]{W-Value (Energy per Ion Pair)}{%
The $W$-value of a material is the average energy deposited per electron--ion
pair created:
\begin{equation}
  \label{eq:lec07:Wvalue}
  W = \frac{E_\text{deposited}}{N_\text{ion pairs}}.
\end{equation}
Typical values: $W_\text{gas} \approx 30\,\text{eV}$,
$W_\text{silicon} \approx 3\,\text{eV}$.}

\paragraph{Where the projectile mass does appear.}
The formula $-dE/dx \propto z^2/v^2$ is a function of \emph{velocity}.  When
we express it as a function of kinetic energy $E_k = \frac{1}{2}Mv^2$
(non-relativistically), we get $v^2 = 2E_k/M$, so
\begin{equation}
  \label{eq:lec07:dEdx_energy}
  -\frac{dE}{dx} \propto \frac{z^2 M}{E_k}.
\end{equation}
The projectile mass enters \emph{only through the velocity-energy conversion}.
At the \emph{same velocity} the formula is independent of $M$;
at the \emph{same kinetic energy} it scales as $z^2 M$.

\paragraph{Comparing $\alpha$ and proton at the same kinetic energy.}
Take $E_\alpha = E_p$.  Since $m_\alpha = 4m_p$,
\begin{equation}
  \label{eq:lec07:v_ratio}
  \frac{1}{2}m_\alpha v_\alpha^2 = \frac{1}{2}m_p v_p^2
  \;\Rightarrow\;
  \frac{v_\alpha^2}{v_p^2} = \frac{m_p}{m_\alpha} = \frac{1}{4}.
\end{equation}
With $z_\alpha = 2$ and $z_p = 1$,
\begin{equation}
  \label{eq:lec07:stop_ratio}
  \frac{(-dE/dx)_\alpha}{(-dE/dx)_p}
  = \frac{z_\alpha^2}{z_p^2}\cdot\frac{v_p^2}{v_\alpha^2}
  = 4 \times 4 = \boxed{16}.
\end{equation}
\textit{Significance:} An $\alpha$ particle loses energy \emph{sixteen times faster}
than a proton of the same kinetic energy.  Heavier and more highly-charged
projectiles have shorter ranges.

\paragraph{Why electrons have larger ranges than protons at the same energy.}
An electron at energy $E_k$ has $v^2 = 2E_k/m_e$, which is $\approx 1836\times$
larger than a proton's $v^2$ at the same energy.  Since $-dE/dx \propto 1/v^2$,
the stopping power on the electron is $\sim 1836\times$ smaller.  (At typical
decay energies electrons are already ultra-relativistic, making the contrast
even sharper.)

%%───────────────────────────────────────────────────────────────────
\subsection{Detector Types — Comparison}
\label{subsec:lec07:detector_comparison}

\subsubsection*{Conceptual Background}
Three detector families were introduced at the end of Lecture~6.  Here we
compare them systematically across the properties that matter experimentally:
energy resolution, timing resolution, cost, and practical constraints.

\begin{table}[h]
\centering
\small
\begin{tabular}{lccc}
\toprule
Property & Scintillator & Gas detector & Solid-state \\
\midrule
Energy resolution & moderate ($\sim$few\,\%) & poor–moderate ($\sim$few\,\%) & excellent ($\sim 0.1\,\%$) \\
Timing resolution & excellent ($\sim$ps–ns) & poor ($\sim\mu$s–ms) & good ($\sim$ns) \\
Cost per area & low & moderate & very high \\
Scalable to large size & yes (cheap plastic) & yes (gas) & no \\
Works as reaction target & no & yes & no \\
\bottomrule
\end{tabular}
\caption{Comparison of the three main detector families.}
\label{tab:lec07:detectors}
\end{table}

\paragraph{Why gas detectors are slow.}
Drift velocity of electrons in gas is $\sim 1\,\text{cm}/\mu\text{s}$,
so a detector of size $\sim\text{cm}$ takes $\sim\mu\text{s}$ to collect
the signal.  This limits the maximum counting rate to $\sim 10$–$100\,\text{kHz}$.

\paragraph{Advantages of gas detectors despite low speed.}
Gas can \emph{surround} a nuclear reaction, yielding $4\pi$ solid-angle
coverage and near-unity geometric efficiency.  Reactions can occur
\emph{inside} the detector volume.  This is the main reason gas detectors
remain in use despite their slow response.

%%───────────────────────────────────────────────────────────────────
\subsection{Statistical Origin of Energy Resolution}
\label{subsec:lec07:resolution}

\subsubsection*{Mathematical Setup}
All three detector types ultimately convert deposited energy into a number of
charge carriers (electrons or photon-electron pairs).  The statistical
fluctuation in this number sets a fundamental lower bound on energy resolution.

\paragraph{Number of charge carriers.}
A particle depositing energy $E$ produces, on average,
\begin{equation}
  \label{eq:lec07:N_carriers}
  \boxed{\bar{N} = \frac{E}{W},}
\end{equation}
where $W$ is the $W$-value.  For Poisson statistics (independent events),
\begin{equation}
  \label{eq:lec07:sigma_N}
  \sigma_N \approx \sqrt{\bar{N}}.
\end{equation}

\paragraph{Relative energy resolution.}
Since $E \propto N$, the relative energy resolution is
\begin{equation}
  \label{eq:lec07:energy_res}
  \boxed{\frac{\sigma_E}{E} = \frac{\sigma_N}{\bar{N}} = \frac{1}{\sqrt{\bar{N}}} = \sqrt{\frac{W}{E}}.}
\end{equation}
This is the \emph{ideal} (quantum-limited) resolution.

\ex[ex:resolution_compare]{Resolution for a 13 MeV Alpha Particle}{%
For $E_\alpha = 13\,\text{MeV}$:

\textbf{Gas detector} ($W \approx 30\,\text{eV}$):
\begin{equation*}
  \bar{N}_\text{gas} \approx \frac{13\times10^6}{30} \approx 4\times10^5,
  \qquad
  \frac{\sigma_E}{E} \approx \frac{1}{\sqrt{4\times10^5}} \approx 0.5\,\%.
\end{equation*}

\textbf{Solid-state} ($W \approx 3\,\text{eV}$):
\begin{equation*}
  \bar{N}_\text{solid} \approx \frac{13\times10^6}{3} \approx 4\times10^6,
  \qquad
  \frac{\sigma_E}{E} \approx \frac{1}{\sqrt{4\times10^6}} \approx 0.1\,\%.
\end{equation*}

The factor $\sim 3$ improvement in resolution comes from the factor
$\sim 10$ larger number of charge carriers, i.e., $\sqrt{10}\approx 3.16$.}

\textit{Significance:} In practice, gas detectors achieve resolutions
significantly worse than this ideal (electronic noise, recombination, gain
non-uniformity), while solid-state detectors routinely approach the ideal
$\sim 0.1\,\%$, making them the instrument of choice for precision spectroscopy.

\nt{For scintillators the analogous quantity is the light yield: the number of
optical photons produced per MeV.  The photon-to-electron conversion in the PMT
introduces additional fluctuations; typical energy resolution is a few percent.}

%%───────────────────────────────────────────────────────────────────
\subsection{Detector Efficiency}
\label{subsec:lec07:efficiency}

\subsubsection*{Physical Motivation}
A detector is only useful if it actually registers the particles we want to
measure.  The \emph{efficiency} tells us what fraction of emitted particles
ends up contributing to our data.  For rare processes — neutrinos, dark matter
candidates, exotic nuclear reactions — maximising efficiency is as important as
any other design consideration.

\dfn[dfn:efficiency]{Detector Efficiency}{%
\begin{equation}
  \label{eq:lec07:efficiency}
  \varepsilon \equiv \frac{N_\text{detected}}{N_\text{emitted}}.
\end{equation}
This quantity depends on the specific definition of ``emitted,'' which must
always be stated clearly when quoting an efficiency.}

\paragraph{Two components of efficiency.}

\textbf{Geometric acceptance} ($\Omega/4\pi$): A point source emitting
isotropically at distance $d$ from a detector of area $A$ subtends a solid
angle $\Omega \approx A/d^2$ (for $d \gg \sqrt{A}$).  The fraction of particles
that reach the detector is
\begin{equation}
  \label{eq:lec07:geo_eff}
  \varepsilon_\text{geo} = \frac{\Omega}{4\pi} \approx \frac{A}{4\pi d^2}.
\end{equation}
This scales as $1/d^2$ — moving the detector twice as far away reduces geometric
efficiency by a factor of 4.

\textbf{Intrinsic efficiency}: Of the particles that enter the detector, what
fraction actually produce a signal?
\begin{itemize}
  \item Charged particles traversing an active volume: $\varepsilon_\text{int} \approx 1$
  (every ionising particle deposits energy).
  \item Photons: binary interaction, exponential probability.  A photon of
  attenuation coefficient $\mu$ in a detector of thickness $x$ has
  \begin{equation}
    \label{eq:lec07:photon_int_eff}
    \varepsilon_\text{int} = 1 - e^{-\mu x} \approx \mu x \quad (\mu x \ll 1).
  \end{equation}
  To detect photons efficiently, one wants large $\mu$ (high-$Z$ material) and
  large $x$, but larger $x$ costs spatial resolution and money.
\end{itemize}

\paragraph{Total efficiency and its importance.}
\begin{equation}
  \label{eq:lec07:total_eff}
  \boxed{\varepsilon_\text{tot} = \varepsilon_\text{geo} \times \varepsilon_\text{int}.}
\end{equation}
The statistical uncertainty on a rate measurement is $\propto 1/\sqrt{N_\text{detected}}
\propto 1/\sqrt{\varepsilon_\text{tot}}$.  Doubling efficiency is equivalent to quadrupling
measurement time.

\nt{DUNE (neutrino experiment): $6\,\text{m} \times 14\,\text{m} \times 62\,\text{m}$
detector volume containing 17\,000 tonnes of liquid argon at cryogenic
temperatures.  The extreme scale is necessary because the cross section for
neutrino-nucleus interactions is $\sim 10^{-38}\,\text{cm}^2$, requiring
enormous target mass to achieve measurable statistics.}

\paragraph{Strategies for maximising efficiency.}
\begin{enumerate}
  \item Enlarge detector area or use an array of detectors.
  \item Move detector closer to source (geometric efficiency $\propto 1/d^2$).
  \item Surround the reaction (gas detectors can provide $4\pi$ coverage).
  \item Improve data acquisition: fast electronics can reduce dead time and
  allow a higher fraction of events to be recorded.
  \item Improve event selection: using trigger criteria to reduce background
  can effectively raise the signal fraction even if the raw geometric efficiency
  stays the same.
\end{enumerate}

%%───────────────────────────────────────────────────────────────────
\subsection{Particle Tracking}
\label{subsec:lec07:tracking}

\subsubsection*{Physical Motivation}
Energy and timing alone often do not tell us everything.  Knowing the
\emph{trajectory} of a particle provides:
its direction of emission (→ reaction kinematics), whether it scattered or
decayed in flight (→ topology), and its curvature in a magnetic field
(→ momentum and charge sign).

\paragraph{General principle.}
A tracker measures the position of a particle at several known planes or
volumes.  By fitting a curve to these space points, one reconstructs the
3D trajectory.

\subsubsection*{Multi-Wire Proportional Chamber (MWPC)}

A MWPC is a gas detector in which the usual single anode wire is replaced by
an array of many parallel wires with pitch $\sim 1$–$2\,\text{mm}$.  Each wire
is read out independently.

\begin{figure}[h]
\centering
\begin{tikzpicture}[>=Stealth, thick, scale=0.9]
  % cathode planes
  \fill[gray!15] (-3,-1.3) rectangle (3,-1.0);
  \fill[gray!15] (-3,+1.0) rectangle (3,+1.3);
  \node[right] at (3.1,1.15) {\small cathode};
  \node[right] at (3.1,-1.15) {\small cathode};
  % anode wires
  \foreach \x in {-2.5,-1.5,-0.5,0.5,1.5,2.5}
    \draw[blue!70!black, line width=1.2pt] (\x,-0.9) -- (\x,+0.9);
  \node[right] at (3.1,0) {\small anode wires};
  % field lines (schematic)
  \draw[orange!60, dashed, <-] (-2.9,0.7) .. controls (-2.9,0.1) .. (-2.5,0);
  \draw[orange!60, dashed, <-] (0.1,0.7) .. controls (0.1,0.1) .. (0.5,0);
  % particle track (arrow)
  \draw[->, red!70!black, very thick] (-2.8,1.5) -- (1.2,-1.5)
    node[right] {\small track};
  % ionisation dots
  \filldraw[red] (-1.8,0.9) circle (2pt);
  \filldraw[red] (-0.9,0.3) circle (2pt);
  \filldraw[red] (0.0,-0.3) circle (2pt);
\end{tikzpicture}
\caption{Schematic MWPC cross-section.  A charged particle crossing the active
volume produces ion clusters (red dots).  Electrons drift to the nearest
anode wire and trigger avalanche multiplication.  The wire that fires identifies
the $x$-coordinate of the hit with $\sim$mm precision.  Two chambers oriented
perpendicularly provide 2D (and, combining with drift time, 3D) information.}
\label{fig:lec07:mwpc}
\end{figure}

\textbf{How the MWPC gives position.}
Near each wire the electric field $\propto 1/r$ causes avalanche multiplication
close to the wire surface.  The avalanche signal is confined to the wire that
was closest to the ionisation cluster.  Reading out each wire separately maps
the \emph{hit position} to a single wire channel — i.e., to a known $x$-coordinate.

A second MWPC with wires rotated by $90^\circ$ provides $y$-coordinates.
Combining both planes gives a 2D hit position; adding the drift time in the gas
gives a third coordinate.

\textbf{Limitations and successors.}
MWPC wire pitch is limited to $\sim 1\,\text{mm}$; stringing thousands of wires
is labour-intensive.  Modern replacements are \emph{Micro-Pattern Gas Detectors}
(MPGD):
\begin{itemize}
  \item \textbf{Micromegas}: fine metallic mesh $\sim 100\,\mu\text{m}$ above a
  patterned electrode; strong field below mesh produces avalanche.
  \item \textbf{GEM (Gas Electron Multiplier)}: insulating foil with micro-holes;
  each hole acts as an independent avalanche region.
\end{itemize}
Both achieve spatial resolution $\sim 50$–$100\,\mu\text{m}$ and can be
fabricated by photolithographic techniques.

\subsubsection*{Silicon Strip and Pixel Detectors}

A thin ($\lesssim 1\,\text{mm}$) silicon diode is segmented into narrow
strips (pitch $\sim 50$–$100\,\mu\text{m}$) or pixels.  Each strip/pixel
is read out independently.  A charged particle crossing the detector creates
electron-hole pairs; the strip closest to the track fires.

By stacking several layers (each with orthogonal strip orientation) one
reconstructs a full 3D track, as illustrated in Fig.~\ref{fig:lec07:si_tracker}.

\begin{figure}[h]
\centering
\begin{tikzpicture}[>=Stealth]
  % Layer 1
  \fill[blue!10] (-2.5,-0.05) rectangle (2.5,0.05);
  \draw[blue!60, thin] (-2.5,0) -- (2.5,0);
  \foreach \x in {-2,-1,0,1,2}
    \draw[blue!30, very thin] (\x,-0.05) -- (\x,0.05);
  % Layer 2
  \fill[blue!10] (-2.5,0.75) rectangle (2.5,0.85);
  \draw[blue!60, thin] (-2.5,0.8) -- (2.5,0.8);
  \foreach \x in {-2,-1,0,1,2}
    \draw[blue!30, very thin] (\x,0.75) -- (\x,0.85);
  % Layer 3
  \fill[blue!10] (-2.5,1.55) rectangle (2.5,1.65);
  \draw[blue!60, thin] (-2.5,1.6) -- (2.5,1.6);
  \foreach \x in {-2,-1,0,1,2}
    \draw[blue!30, very thin] (\x,1.55) -- (\x,1.65);
  % Layer 4
  \fill[blue!10] (-2.5,2.35) rectangle (2.5,2.45);
  \draw[blue!60, thin] (-2.5,2.4) -- (2.5,2.4);
  \foreach \x in {-2,-1,0,1,2}
    \draw[blue!30, very thin] (\x,2.35) -- (\x,2.45);
  % particle track
  \draw[->, red!80!black, thick]
    (-2.2,-0.3) -- (1.8,2.7) node[right] {\small track};
  % hit marks (explicit)
  \draw[red, thick] (-1.85,0) +(-.1,-.1) -- +(.1,.1)
                    (-1.85,0) +(-.1,.1)  -- +(.1,-.1);
  \draw[red, thick] (-0.9,0.8) +(-.1,-.1) -- +(.1,.1)
                    (-0.9,0.8) +(-.1,.1)  -- +(.1,-.1);
  \draw[red, thick] (0.05,1.6) +(-.1,-.1) -- +(.1,.1)
                    (0.05,1.6) +(-.1,.1)  -- +(.1,-.1);
  \draw[red, thick] (1.0,2.4) +(-.1,-.1) -- +(.1,.1)
                    (1.0,2.4) +(-.1,.1)  -- +(.1,-.1);
  % labels
  \node[right] at (2.6,0) {\small layer 1};
  \node[right] at (2.6,0.8) {\small layer 2};
  \node[right] at (2.6,1.6) {\small layer 3};
  \node[right] at (2.6,2.4) {\small layer 4};
\end{tikzpicture}
\caption{Four silicon strip detector layers.  Each layer records a hit
position (cross marks).  Connecting the hits reconstructs the 3D trajectory.
Alternating strip orientations ($0^\circ$, $90^\circ$, $\ldots$) in successive
layers provide both transverse coordinates.}
\label{fig:lec07:si_tracker}
\end{figure}

\nt{This is the principle of the central tracking detectors in ATLAS and CMS at
the LHC.  The ``giant wheel'' surrounding the beam axis is largely a silicon
pixel tracker plus outer strip layers, whose primary purpose is to reconstruct
particle trajectories and measure their momenta via curvature in the axial
magnetic field.}

%%───────────────────────────────────────────────────────────────────
\subsection{Particle Identification}
\label{subsec:lec07:pid}

\subsubsection*{Physical Motivation}
So far detectors measured only energy; they were agnostic to particle species.
In nuclear physics experiments one often needs to know \emph{which particle}
produced the signal (proton, deuteron, $\alpha$, heavier fragment?).
Several techniques exploit the mass and charge dependence of energy loss.

\subsubsection*{Method 1: $dE/dx$–Range}

From the Bethe-Bloch formula, at the same energy $E_k$,
$-dE/dx \propto z^2 M / E_k$ and the range $R \propto E_k/(z^2 M)$.
A gas (or silicon) detector that is thick enough to stop the particle
records both the total deposited energy \emph{and} the path length
(from the track or the stopping position).  The combination
$(dE/dx,\, \text{range})$ uniquely identifies $(z, M)$.

\subsubsection*{Method 2: $\Delta E$–$E$ Telescope}

\dfn[dfn:dEE_telescope]{$\Delta E$--$E$ Telescope}{%
A $\Delta E$–$E$ telescope consists of two detectors in series:
\begin{enumerate}
  \item A \textbf{thin} ($\Delta E$) detector: the particle passes through,
  depositing a small fraction of its energy.
  \item A \textbf{thick} ($E$) detector: the particle stops here, depositing
  its remaining energy.
\end{enumerate}
One records the pair $(\Delta E,\, E_\text{stop})$ for each event.}

\textit{Why does this identify the particle?}
For a non-relativistic particle of charge $z$, mass $M$, and kinetic energy $E_k$,
\begin{equation}
  \label{eq:lec07:dEdx_telescope}
  -\frac{dE}{dx} \propto \frac{z^2 M}{E_k}.
\end{equation}
The energy lost in the $\Delta E$ layer is approximately
$\Delta E \approx (dE/dx)\cdot\Delta x \propto z^2 M / E_k$.
But $E_k = \Delta E + E_\text{stop}$, so
\begin{equation}
  \label{eq:lec07:dEE_relation}
  \Delta E \cdot (\Delta E + E_\text{stop}) \propto z^2 M.
\end{equation}
Plotting $\Delta E$ vs.\ $E_\text{stop}$, each species $(z, M)$ traces out a
\textbf{characteristic hyperbola-like curve}.  Different isotopes and elements
separate cleanly on this plot.

\begin{figure}[h]
\centering
\begin{tikzpicture}[>=Stealth]
  \draw[->] (0,0) -- (5.5,0) node[right] {$E_\text{stop}$};
  \draw[->] (0,0) -- (0,3.5) node[above] {$\Delta E$};
  % proton curve
  \draw[blue!70!black, thick, smooth]
    plot[tension=0.7] coordinates {(0.3,2.8)(0.8,1.8)(1.8,1.0)(3.5,0.5)(5.0,0.3)};
  \node[blue!70!black] at (1.2,2.3) {\small $p$};
  % deuteron curve
  \draw[green!60!black, thick, smooth]
    plot[tension=0.7] coordinates {(0.4,3.1)(1.0,2.2)(2.2,1.3)(4.0,0.7)};
  \node[green!60!black] at (1.6,2.7) {\small $d$};
  % alpha curve
  \draw[red!70!black, thick, smooth]
    plot[tension=0.7] coordinates {(0.2,3.3)(0.5,2.8)(1.0,2.2)(2.0,1.6)(3.5,1.1)};
  \node[red!70!black] at (0.7,3.15) {\small $\alpha$};
\end{tikzpicture}
\caption{Schematic $\Delta E$–$E$ plot.  Each particle species traces a
distinct locus.  Protons ($z{=}1$, $M{=}1$), deuterons ($z{=}1$, $M{=}2$),
and $\alpha$ particles ($z{=}2$, $M{=}4$) separate cleanly, allowing
unambiguous identification.}
\label{fig:lec07:dEE}
\end{figure}

\textit{Significance:} The $\Delta E$–$E$ telescope is one of the most widely
used particle identification methods in nuclear physics.  It requires only two
detectors and is straightforward to implement with either scintillators or
solid-state devices.

\subsubsection*{Method 3: Gamma–Charged Coincidence}

The $\gamma$-ray spectrum emitted by a nucleus is a unique fingerprint of that
isotope (analogous to atomic emission lines).  If a nuclear reaction produces a
specific isotope in an excited state, the characteristic de-excitation $\gamma$
rays can be measured in coincidence with the charged particle.  Together,
$(E_\text{charged},\, E_\gamma)$ identify both the isotope and its excitation energy.

%%───────────────────────────────────────────────────────────────────
\subsection{Time-of-Flight for Neutron Energy Measurement}
\label{subsec:lec07:tof}

\subsubsection*{Physical Motivation}
Neutrons have no charge and cannot be detected directly.  The standard approach
is \emph{conversion}: a nuclear reaction produces a charged particle that can
be detected.  However, for $^6\text{Li}$-based converters (reaction
$n + {}^6\text{Li} \to {}^4\text{He} + {}^3\text{H}$), the light output is
dominated by the $Q$-value of the conversion reaction ($Q = 4.78\,\text{MeV}$),
\emph{not} by the incoming neutron energy.  There is therefore no direct
energy-to-light correlation, and one cannot measure $E_n$ from pulse height
alone.

\dfn[dfn:tof]{Time-of-Flight (ToF) Method}{%
Produce neutrons in \textbf{pulses} at a known position.  Start a clock when
the pulse is produced (e.g., using the beam pick-up signal from the accelerator);
stop the clock when a neutron is detected at distance $d$.  The neutron energy
is then:
\begin{equation}
  \label{eq:lec07:tof}
  \boxed{E_n = \frac{1}{2}m_n\left(\frac{d}{t_\text{ToF}}\right)^2,}
\end{equation}
valid for $E_n \ll m_n c^2 = 939.6\,\text{MeV}$ (non-relativistic regime).}

\textit{Significance:} ToF is the standard technique for measuring neutron
energy spectra in nuclear physics.  The energy resolution improves with longer
flight paths (better time-to-energy discrimination) but at the cost of reduced
geometric efficiency.

\paragraph{Experimental example.}
In a recent experiment at a small accelerator in Braunschweig (Germany), protons
impinged on an ${}^{18}\text{O}$ target:
\begin{equation}
  p + {}^{18}\text{O} \to {}^{18}\text{F} + n.
\end{equation}
The neutron energy spectrum was measured by ToF using a $^6\text{LiF}$ scintillator
detector.  The proton beam was pulsed; the beam pick-up signal started the clock,
and the scintillator signal stopped it.  From Eq.~\eqref{eq:lec07:tof}, each
recorded time difference converts directly to a neutron energy.

%%───────────────────────────────────────────────────────────────────
\subsection{Experimental Spectra}
\label{subsec:lec07:spectra}

\subsubsection*{Alpha Spectrum of the $^{228}$Th Decay Chain}

\subsubsection*{Physical Motivation}
Thorium-228 undergoes a long series of successive $\alpha$ and $\beta$ decays.
Each $\alpha$-decaying isotope in the chain emits $\alpha$ particles at a
\emph{unique, discrete} energy — a direct consequence of energy-momentum
conservation and the discrete nuclear energy levels.

The chain ${}^{228}\text{Th} \to \cdots \to {}^{208}\text{Pb}$ includes:
\begin{equation*}
  {}^{228}\text{Th} \xrightarrow{\alpha}
  {}^{224}\text{Ra} \xrightarrow{\alpha}
  {}^{220}\text{Rn} \xrightarrow{\alpha}
  {}^{216}\text{Po} \xrightarrow{\alpha}
  {}^{212}\text{Pb} \xrightarrow{\beta}
  {}^{212}\text{Bi} \xrightarrow{\beta+\alpha}
  {}^{208}\text{Pb}.
\end{equation*}
A silicon detector placed near a $^{228}$Th source records all $\alpha$ peaks
simultaneously, since the chain quickly reaches secular equilibrium.

\textbf{Key features of the measured spectrum:}
\begin{itemize}
  \item \textbf{Sharp peaks}: each peak corresponds to one $\alpha$-emitting
  isotope; peak width reflects the detector energy resolution ($\sim 0.1\,\%$
  for silicon at room temperature).
  \item \textbf{No background}: charged particles deposit \emph{all} their
  energy in the detector; there is no Compton continuum analogue.
  \item \textbf{Low-energy tail}: ``incomplete energy collection'' — some
  $\alpha$-produced ions recombine near the surface or in the dead layer before
  being swept out, leading to a slight tail below the main peak.
  This is a detector artefact, not physics.
  \item \textbf{Axis in channels}: the ADC converts the signal height (voltage)
  to a channel number.  Calibration maps channel $\leftrightarrow$ energy using
  the known $\alpha$ energies of the isotopes in the chain.
\end{itemize}

\subsubsection*{Gamma Spectrum of $^{60}$Co with HPGe Detector}

$^{60}$Co decays via $\beta^-$ to an excited state of $^{60}$Ni, which
de-excites by emitting two $\gamma$ rays in cascade:
$E_{\gamma_1} \approx 1173\,\text{keV}$ and $E_{\gamma_2} \approx 1333\,\text{keV}$.

An HPGe (high-purity germanium) detector provides energy resolution of
$\sim 0.1\,\%$, resolving these two lines clearly.  The spectrum shows
(on a log scale):

\paragraph{Photo peaks.}
Two narrow peaks at $1173$ and $1333\,\text{keV}$.  A photon is completely
absorbed by the photoelectric effect (or a sequence of interactions whose net
result is complete energy deposition); all its energy goes to one electron
(minus the small binding energy).  The resulting sharp spike is the
\emph{full-energy peak} (commonly called the \emph{photo peak}).

\paragraph{Compton continuum.}
The dominant feature across most of the spectrum.  A $\gamma$ scatters
off a quasi-free electron (Compton scattering), depositing a \emph{continuous}
range of electron energies from near zero up to the Compton edge:
\begin{equation}
  \label{eq:lec07:compton_edge_spectrum}
  T_{e,\max} = E_\gamma - \frac{E_\gamma}{1 + 2E_\gamma/m_e c^2}.
\end{equation}
For $E_{\gamma_1} = 1173\,\text{keV}$, this gives $T_{e,\max} \approx 963\,\text{keV}$;
for $E_{\gamma_2} = 1333\,\text{keV}$, $T_{e,\max} \approx 1118\,\text{keV}$.
The spectrum shows \textbf{two Compton edges} — characteristic step-like drops
at these energies — visible even on the log scale.

\paragraph{Log scale is essential.}
On a linear scale only the two photo peaks are visible; all other features are
invisible.  The log scale reveals:
\begin{enumerate}
  \item Two Compton edges (as sharp drops in the continuum).
  \item The region between the Compton edge and the photo peak
  (``Compton valley'') contains contributions from multiple scatters.
  \item A $511\,\text{keV}$ annihilation peak (weak, from pair production
  followed by positron annihilation) is visible in principle but requires
  a strong source.
\end{enumerate}

\nt{The photo peak is not \emph{exclusively} from the photoelectric effect.
Compton scattering followed by complete absorption of the scattered photon also
deposits the full photon energy.  Monte Carlo simulations show that pair
production (above threshold) and multiple Compton scatters all contribute to
the full-energy peak.}

%%───────────────────────────────────────────────────────────────────
\subsection*{Key Takeaways}
%%───────────────────────────────────────────────────────────────────
\begin{itemize}
  \item \textbf{Bethe-Bloch correction:} the projectile mass $M$ does
  not appear in $-dE/dx(v)$; it enters only when converting to
  $-dE/dx(E_k)$ via $v^2 = 2E_k/M$.  At the same energy, stopping power
  ratios go as $z^2 M$.  An $\alpha$ is stopped $16\times$ more efficiently
  than a proton at the same $E_k$.
  \item \textbf{Energy resolution:} limited by $\sigma_E/E = \sqrt{W/E}$.
  Silicon ($W \approx 3\,\text{eV}$) gives $\sim 0.1\,\%$; gas ($W \approx
  30\,\text{eV}$) gives ideal $\sim 0.5\,\%$ but worse in practice.
  \item \textbf{Efficiency:} $\varepsilon_\text{tot} = \varepsilon_\text{geo}
  \times \varepsilon_\text{int}$; geometric efficiency $\propto 1/d^2$.
  For rare processes, maximising efficiency is a primary design driver.
  \item \textbf{Tracking:} MWPC, MPGD (micromegas/GEM), silicon strip/pixel
  detectors all reconstruct 3D particle trajectories; silicon gives
  $\sim 50\,\mu\text{m}$ resolution.
  \item \textbf{Particle ID:} $\Delta E$–$E$ telescope identifies species via
  $\Delta E \cdot E_\text{tot} \propto z^2 M$; each species maps to a distinct
  curve in the $(\Delta E, E_\text{stop})$ plane.
  \item \textbf{Neutron energy:} measured by ToF since direct energy-to-light
  conversion fails in conversion detectors.
  $E_n = \tfrac{1}{2}m_n(d/t)^2$ (non-relativistic).
\end{itemize}
